% H1/G3-bis: ぶら下げ（hanging_punctuation, -b 変種）の実測検証
% pdfTeX の \pdfprotrudechars + \rpcode による句読点の行末突出し:
%   case A: クラス7（、）, case B: クラス6（。）, case C: 混在長文の安定性。
% 検証（すべてログ内判定）:
%   - 突出し機構の発火: 行ボックスの右マージンカーンを \rightmarginkern で
%     読み戻す（-10.0pt = -1em で固定: rpcode=1000 × quad）。グリフ全幅1zw
%     の突出し量を持つマージンカーンが組み込まれていれば HANG-PASS。
%   - 段落行数: 漢×9、漢×9 が hsize 95pt で2行に組めること（\prevgraf=2）。
% 注意: このフィクスチャは段落中に \pdfsavepos 等の whatsit を置かない。
% 、の直後の補正カーンとペアグルーの間に whatsit が入るとブレイク成立条件
% が変わり突出しが発火しないことがある（poc/README-g3b.md 参照）。
\documentclass[pdflatex,hanging_punctuation]{jlreq}
\makeatletter
\parindent=0pt


\begin{document}
% ---- case A: 読点「、」の行末突出し
% hsize=95pt（全角9.5文字分）: 漢×9、漢×9 → 1行目=[漢×9][、+マージンカーン]，
% 2行目=[漢×9] の2行組になるはず。
\setbox10=\vbox{\hsize=95pt
  漢漢漢漢漢漢漢漢漢、漢漢漢漢漢漢漢漢漢\par
  \global\setbox11=\lastbox \unskip\unpenalty
  \global\setbox12=\lastbox
  \global\count1=\prevgraf}
\unvbox10
\edef\h@rmk{\rightmarginkern12}
\ifnum\pdfstrcmp{\h@rmk}{-10.0pt}=0 
  \typeout{HANG-PASS:A margin-kern=\h@rmk}%
\else
  \typeout{HANG-FAIL:A margin-kern=[\h@rmk] expected=[-10.0pt]}%
\fi
\ifnum\count1=2
  \typeout{HANG-PASS:A lines=2}%
\else
  \typeout{HANG-FAIL:A lines=\the\count1 expected=2}%
\fi

% ---- case B: 句点「。」の行末突出し（クラス6）
\setbox10=\vbox{\hsize=95pt
  漢漢漢漢漢漢漢漢漢。漢漢漢漢漢漢漢漢漢\par
  \global\setbox11=\lastbox \unskip\unpenalty
  \global\setbox12=\lastbox
  \global\count1=\prevgraf}
\unvbox10
\edef\h@rmk{\rightmarginkern12}
\ifnum\pdfstrcmp{\h@rmk}{-10.0pt}=0 
  \typeout{HANG-PASS:B margin-kern=\h@rmk}%
\else
  \typeout{HANG-FAIL:B margin-kern=[\h@rmk] expected=[-10.0pt]}%
\fi
\ifnum\count1=2
  \typeout{HANG-PASS:B lines=2}%
\else
  \typeout{HANG-FAIL:B lines=\the\count1 expected=2}%
\fi

% ---- case C: 混在長文の安定性（ぶら下げ＋禁則＋ペアグルーが共存して
% エラーなく組めること：行数だけ簡易確認）
\setbox10=\vbox{\hsize=200pt
  昔々，あるところにお爺さんとお婆さんがありました。お爺さんは山へ
  芝刈りに，お婆さんは川へ洗濯に行きました。川を下ってきた大きな桃を
  拾って帰ると，桃が割れて中から元気な男の子が飛び出しました。%
  \par\global\count1=\prevgraf}
\unvbox10
\ifnum\count1>0
  \typeout{HANG-PASS:C lines=\the\count1}%
\else
  \typeout{HANG-FAIL:C lines=\the\count1}%
\fi

\end{document}
